% GENERATED FILE. Edit canonical lessons, then run npm run generate:books.
% canonical-source-hash: fnv1a64:e5364e08542ab9dd
% canonical-lessons: FR-C345-approuver, FR-C345-desapprouver, FR-C345-partisan, FR-C345-adversaire, FR-C345-argumenter

\chapter{Pour ou contre}
\label{ch:fr-pour-ou-contre}
\hlchaptermodality{345}

\begin{chapteropening}
\textbf{By the end of this chapter:} \emph{I can say who is for or against an idea: approve or disapprove, name a supporter and an opponent, and argue a case.}

Complete the last lesson of chapter 345: I can say who is for or against an idea: approve or disapprove, name a supporter and an opponent, and argue a case.
\end{chapteropening}

\section[approuver]{approuver — to approve (of)}

\label{lesson:FR-C345-approuver}



\begin{quote}
Before the new one: say the French for the subtle difference, then the French for to qualify.
\end{quote}

\subsection*{You'll want to know: approuver}
\textbf{approuver} — \textquotedblleft{}to approve (of)\textquotedblright{}.

\begin{quote}
Mes parents n'approuvent pas mon choix.
\end{quote}

\begin{quote}
\emph{My parents do not approve of my choice.}
\end{quote}

\emph{Approuver une décision} — to approve of a decision.

\begin{cousinweb}
\emph{approuver} comes from Latin \emph{approbāre} \textquotedblleft{}to approve, to assent to\textquotedblright{} (EtymDB), built from \emph{ad} and \emph{probāre} \textquotedblleft{}to test, to find good\textquotedblright{} (Lewis \& Short).
\end{cousinweb}

\begin{grammarlens}[title={What you've built}]
One more word for telling, explaining and describing at B1.
\end{grammarlens}

\subsection*{Your turn}
\begin{itemize}
\raggedright
  \item \emph{Say it:} \emph{approuver}
  \item \emph{Say it:} \emph{approuver}, once more
  \item \emph{Say it:} a sentence of your own with \emph{approuver}
  \item \emph{Recall:} say the French for the subtle difference, then the French for to qualify, then say \emph{approuver} again
\end{itemize}

\begin{tcolorbox}[breakable,colback=teal!4,colframe=teal!35!black,title={Before you move on}]
What does approuver mean? (\textquotedblleft{}To approve (of)\textquotedblright{}.) Say it once more.
\end{tcolorbox}
\section[désapprouver]{désapprouver — to disapprove (of)}

\label{lesson:FR-C345-desapprouver}



\begin{quote}
Before the new one: say the French for to qualify, then the French for to approve.
\end{quote}

\subsection*{You'll want to know: désapprouver}
\textbf{désapprouver} — \textquotedblleft{}to disapprove (of)\textquotedblright{}.

\begin{quote}
Le maire désapprouve ce projet.
\end{quote}

\begin{quote}
\emph{The mayor disapproves of this project.}
\end{quote}

\emph{Dés-} + \emph{approuver}: the opposite.

\begin{cousinweb}
\emph{désapprouver} is built from \emph{dés-} \textquotedblleft{}not\textquotedblright{} and \emph{approuver}, which comes from Latin \emph{approbāre} \textquotedblleft{}to approve\textquotedblright{}, on \emph{probāre} \textquotedblleft{}to test\textquotedblright{} (Lewis \& Short).
\end{cousinweb}

\begin{grammarlens}[title={What you've built}]
One more word for telling, explaining and describing at B1.
\end{grammarlens}

\subsection*{Your turn}
\begin{itemize}
\raggedright
  \item \emph{Say it:} \emph{désapprouver}
  \item \emph{Say it:} \emph{désapprouver}, once more
  \item \emph{Say it:} a sentence of your own with \emph{désapprouver}
  \item \emph{Recall:} say the French for to qualify, then the French for to approve, then say \emph{désapprouver} again
\end{itemize}

\begin{tcolorbox}[breakable,colback=teal!4,colframe=teal!35!black,title={Before you move on}]
What does désapprouver mean? (\textquotedblleft{}To disapprove (of)\textquotedblright{}.) Say it once more.
\end{tcolorbox}
\section[le partisan]{le partisan — the supporter (of an idea)}

\label{lesson:FR-C345-partisan}



\begin{quote}
Before the new one: say the French for to approve, then the French for to disapprove.
\end{quote}

\subsection*{You'll want to know: le partisan}
\textbf{le partisan} — \textquotedblleft{}the supporter (of an idea)\textquotedblright{}.

\begin{quote}
Les partisans du projet sont contents.
\end{quote}

\begin{quote}
\emph{The supporters of the project are happy.}
\end{quote}

\emph{Être partisan de} — to be in favour of. Feminine \emph{partisane}.

\begin{cousinweb}
\emph{partisan} was borrowed from Italian \emph{partigiano} \textquotedblleft{}a supporter of a side\textquotedblright{} (EtymDB), built on \emph{parte} \textquotedblleft{}part, side\textquotedblright{}, from Latin \emph{pars, partis} (Lewis \& Short).
\end{cousinweb}

\begin{grammarlens}[title={What you've built}]
One more word for telling, explaining and describing at B1.
\end{grammarlens}

\subsection*{Your turn}
\begin{itemize}
\raggedright
  \item \emph{Say it:} \emph{le partisan}
  \item \emph{Say it:} \emph{le partisan}, once more
  \item \emph{Say it:} a sentence of your own with \emph{le partisan}
  \item \emph{Recall:} say the French for to approve, then the French for to disapprove, then say \emph{le partisan} again
\end{itemize}

\begin{tcolorbox}[breakable,colback=teal!4,colframe=teal!35!black,title={Before you move on}]
What does le partisan mean? (\textquotedblleft{}The supporter (of an idea)\textquotedblright{}.) Say it once more.
\end{tcolorbox}
\section[l'adversaire]{l'adversaire — the opponent}

\label{lesson:FR-C345-adversaire}



\begin{quote}
Before the new one: say the French for to disapprove, then the French for the supporter.
\end{quote}

\subsection*{You'll want to know: l'adversaire}
\textbf{l'adversaire} — \textquotedblleft{}the opponent\textquotedblright{}.

\begin{quote}
Son adversaire a gagné le débat.
\end{quote}

\begin{quote}
\emph{His opponent won the debate.}
\end{quote}

Masculine or feminine: \emph{un} or \emph{une adversaire}.

\begin{cousinweb}
\emph{adversaire} comes from Latin \emph{adversārius} \textquotedblleft{}an opponent\textquotedblright{} (EtymDB; Lewis \& Short), from \emph{adversus} \textquotedblleft{}turned toward, facing\textquotedblright{}, built on \emph{vertere} \textquotedblleft{}to turn\textquotedblright{}.
\end{cousinweb}

\begin{grammarlens}[title={What you've built}]
One more word for telling, explaining and describing at B1.
\end{grammarlens}

\subsection*{Your turn}
\begin{itemize}
\raggedright
  \item \emph{Say it:} \emph{l'adversaire}
  \item \emph{Say it:} \emph{l'adversaire}, once more
  \item \emph{Say it:} a sentence of your own with \emph{l'adversaire}
  \item \emph{Recall:} say the French for to disapprove, then the French for the supporter, then say \emph{l'adversaire} again
\end{itemize}

\begin{tcolorbox}[breakable,colback=teal!4,colframe=teal!35!black,title={Before you move on}]
What does l'adversaire mean? (\textquotedblleft{}The opponent\textquotedblright{}.) Say it once more.
\end{tcolorbox}
\section[argumenter]{argumenter — to argue a case, to give reasons}

\label{lesson:FR-C345-argumenter}



\begin{quote}
Before the new one: say the French for the supporter, then the French for the opponent.
\end{quote}

\subsection*{You'll want to know: argumenter}
\textbf{argumenter} — \textquotedblleft{}to argue a case, to give reasons\textquotedblright{}.

\begin{quote}
Elle argumente avec calme.
\end{quote}

\begin{quote}
\emph{She argues her case calmly.}
\end{quote}

Not quarrelling: \emph{argumenter} is to back a view with reasons; to quarrel is \emph{se disputer}.

\begin{cousinweb}
\emph{argumenter} is a learned borrowing from Latin \emph{argūmentārī} \textquotedblleft{}to bring forward proof\textquotedblright{} (Lewis \& Short), from \emph{argūmentum} \textquotedblleft{}proof\textquotedblright{}, built on \emph{arguere} \textquotedblleft{}to make clear\textquotedblright{}.
\end{cousinweb}

\begin{grammarlens}[title={What you've built}]
One more word for telling, explaining and describing at B1.
\end{grammarlens}

\subsection*{Your turn}
\begin{itemize}
\raggedright
  \item \emph{Say it:} \emph{argumenter}
  \item \emph{Say it:} \emph{argumenter}, once more
  \item \emph{Say it:} a sentence of your own with \emph{argumenter}
  \item \emph{Recall:} say the French for the supporter, then the French for the opponent, then say \emph{argumenter} again
\end{itemize}

\begin{tcolorbox}[breakable,colback=teal!4,colframe=teal!35!black,title={Before you move on}]
What does argumenter mean? (\textquotedblleft{}To argue a case, to give reasons\textquotedblright{}.) Say it once more.
\end{tcolorbox}
